\documentclass[aspectratio=169]{beamer}
\input{header}
\coursecode{JKSSB}

\title{Hacking, Phishing \& Spoofing}
\subtitle{Lesson 43 --- Attacks and Deception}
\author{JKSSB Computer Awareness}
\date{\today}

\begin{document}
\frame{\titlepage}

\begin{frame}{Why Attacks Succeed}
  \begin{itemize}
    \item Many computer attacks do not break the machine; they
      \key{trick the person} using it.
    \item An attacker needs only two things: a \key{weakness} to use and
      a \key{person} who responds without checking.
    \item Weaknesses can be technical (weak password, old software) or
      human (trust, hurry, fear).
    \item This lesson teaches the names, the signs, and the
      \key{safe response} --- not how to attack.
  \end{itemize}
  \begin{block}{The one idea}
    Attacks use deception or a known weakness; defence is verification
    plus good habits.
  \end{block}
\end{frame}

\begin{frame}{Four Words to Keep Straight}
  \begin{itemize}
    \item \key{Vulnerability}: a weakness in software, settings, or
      human habit that can be misused.
    \item \key{Exploit}: a method that takes advantage of a
      vulnerability.
    \item \key{Attack}: the complete harmful act against a person or
      system.
    \item \key{Payload}: the harmful part delivered by the attack, such
      as stolen data or a malicious file.
  \end{itemize}
  \begin{block}{Read it as a chain}
    Vulnerability (weakness) $\rightarrow$ exploit (method)
    $\rightarrow$ attack (act) $\rightarrow$ payload (harm).
  \end{block}
\end{frame}

\begin{frame}{What Hacking Means}
  \begin{itemize}
    \item \key{Hacking} means gaining \key{unauthorised access} to a
      computer, account, or network.
    \item A \key{hacker} is the person who does it; unauthorised access
      is unlawful even when no data is stolen.
    \item An \key{ethical hacker} tests security \key{with written
      permission} to find and fix weaknesses.
    \item \muted{Exam depth: hacking is the broad act; phishing and
      spoofing are specific deceptive methods inside it.}
  \end{itemize}
  \begin{alertblock}{Trap alert}
    Hacking is unauthorised access. Permission separates an ethical
    security test from an attack.
  \end{alertblock}
\end{frame}

\begin{frame}{Password Attacks}
  \begin{itemize}
    \item \key{Guessing}: trying obvious passwords such as name, birth
      year, or \texttt{password123}.
    \item \key{Dictionary attack}: trying every word from a list of
      common passwords automatically.
    \item \key{Brute-force attack}: trying \key{every possible
      combination} until one works; longer passwords resist it.
    \item \key{Shoulder surfing}: watching the keyboard or screen while
      the victim types.
  \end{itemize}
  \begin{block}{Keep the pair straight}
    Dictionary tries a word list; brute force tries every combination.
  \end{block}
\end{frame}

\begin{frame}{Strong Password Habits}
  \begin{itemize}
    \item Use a \key{long} password with mixed letters, digits, and
      symbols; never reuse it across accounts.
    \item Never share a password or an \key{OTP} (One-Time Password)
      with anyone, including a caller who claims to be bank staff.
    \item Lock the screen when leaving the desk; cover the keypad while
      typing a PIN.
    \item Use \key{two-step verification} where offered: password plus a
      code on the phone.
  \end{itemize}
  \begin{exampleblock}{Office desktop example}
    A clerk who uses the same short password for email and office login
    loses both accounts from one guess; a long unique password with
    two-step verification blocks the same attempt.
  \end{exampleblock}
\end{frame}

\begin{frame}{Social Engineering}
  \begin{itemize}
    \item \key{Social engineering} means tricking a person into giving
      access, instead of breaking software.
    \item \key{Pretexting}: inventing a false story (``bank officer'',
      ``IT support'') to extract details.
    \item \key{Baiting}: offering something free (prize, free recharge,
      free software) to trigger a click or download.
    \item Pressure words --- \key{urgent}, \key{blocked}, \key{verify
      now} --- are the standard signal.
  \end{itemize}
  \begin{alertblock}{Safe response}
    Slow down, verify through a known channel, and never disclose a
    password or OTP under pressure.
  \end{alertblock}
\end{frame}

\begin{frame}{Phishing}
  \begin{itemize}
    \item \key{Phishing} is a fake message, mail, or call that
      \key{imitates a trusted sender} to steal passwords, OTPs, or
      money.
    \item The message carries a \key{malicious link} to a look-alike
      login page, or a \key{malicious attachment} to open.
    \item Common channels: email, \key{SMS smishing}, and voice-call
      \key{vishing}.
    \item Signs: urgent tone, spelling errors, unknown sender, and a
      link that does \key{not} match the real website address.
  \end{itemize}
  \begin{block}{Definition}
    Phishing steals credentials or money by imitation plus a link,
    attachment, or call to action.
  \end{block}
\end{frame}

\begin{frame}{Spoofing}
  \begin{itemize}
    \item \key{Spoofing} means \key{faking an identity}: the sender
      address, caller number, or website looks genuine but is forged.
    \item \key{Email spoofing}: the ``From'' address shows a trusted
      name that was never the real sender.
    \item \key{Caller ID spoofing}: the displayed phone number is
      forged to look like a bank or office.
    \item \key{Website (URL) spoofing}: a look-alike address
      (\texttt{bank-verlfy.com} for \texttt{bank-verify.com}) hosts a
      fake login page.
  \end{itemize}
  \begin{block}{What spoofing is}
    A forged identity used as a disguise; phishing often \key{uses}
    spoofing to look trustworthy.
  \end{block}
\end{frame}

\begin{frame}{Phishing vs Spoofing vs Hacking}
  \begin{center}
    \begin{tabular}{lll}
      \toprule
      \textbf{Term} & \textbf{What it is} & \textbf{One-line test} \\
      \midrule
      Hacking & Unauthorised access & Was access taken without permission? \\
      Phishing & Fraud to steal data/money & Was a fake message asking for action? \\
      Spoofing & Forged identity & Was the sender or address faked? \\
      \bottomrule
    \end{tabular}
  \end{center}
  \vspace{0.35cm}
  \begin{alertblock}{Trap alert}
    A phishing mail \key{uses} a spoofed address; spoofing is the mask,
    phishing is the fraud, hacking is the broad act.
  \end{alertblock}
\end{frame}

\begin{frame}{Malicious Links}
  \begin{itemize}
    \item A \key{malicious link} leads to a fake login page or starts an
      unwanted download.
    \item Hover over the link to read the \key{real address} before
      clicking; on a phone, long-press to preview it.
    \item Type the website address yourself instead of tapping a link in
      an unexpected message.
    \item Shortened or misspelled addresses in an urgent message are a
      warning sign.
  \end{itemize}
  \begin{exampleblock}{Office desktop example}
    A mail says the salary slip is ready and points to
    \texttt{offlce-portal.com}; the real portal is
    \texttt{office-portal.com}. Typing the known address avoids the
    fake page.
  \end{exampleblock}
\end{frame}

\begin{frame}{Malicious Attachments}
  \begin{itemize}
    \item A \key{malicious attachment} is a file that harms the computer
      or steals data when opened.
    \item Do not open an unexpected attachment, even from a known name
      --- confirm with the sender first.
    \item Be cautious with files that ask to ``enable macros'' or
      ``enable content'' on opening.
    \item \muted{Exam depth: attachment is the file; link is the
      address; both can carry the payload.}
  \end{itemize}
  \begin{block}{Safe habit}
    Unexpected file $\rightarrow$ stop, verify the sender, delete if
    unverified.
  \end{block}
\end{frame}

\begin{frame}{Safe Habits That Block Most Attacks}
  \begin{itemize}
    \item Verify the sender through a \key{known channel} before acting
      on any urgent request.
    \item Keep software and \key{antivirus} updated; scan external
      drives before opening files.
    \item Back up important files so one incident cannot destroy all
      copies.
    \item Report phishing messages using the mail client's
      \key{report phishing} option instead of forwarding or replying.
  \end{itemize}
\end{frame}

\begin{frame}{If You Are Targeted}
  \begin{itemize}
    \item Do \key{not} enter a password, OTP, or card details on a page
      reached from a suspicious message.
    \item Disconnect doubt: close the page, delete the message, and do
      not call back numbers given in it.
    \item Change the password from the \key{real website}, and inform
      the bank or office through its published helpline.
    \item Keep a record (screenshot, sender address) for the complaint.
  \end{itemize}
  \begin{alertblock}{Order of action}
    Stop $\rightarrow$ verify on the real channel $\rightarrow$ change
    credentials $\rightarrow$ report.
  \end{alertblock}
\end{frame}

\begin{frame}{Trap Alert: Attack Facts to Keep Straight}
  \begin{itemize}
    \item Phishing \key{imitates} a trusted sender to steal data; spoofing
      \key{forges} the identity it wears.
    \item OTP stands for \key{One-Time Password}; no genuine bank or
      office asks for it by call or message.
    \item Dictionary tries a \key{word list}; brute force tries
      \key{every combination}.
    \item Vulnerability is the \key{weakness}; exploit is the
      \key{method}; payload is the \key{harm delivered}.
    \item Ethical hacking needs \key{written permission}; without it,
      access is unauthorised.
  \end{itemize}
\end{frame}

\begin{frame}{Lesson 43: Recall}
  \begin{itemize}
    \item Hacking is unauthorised access; ethical hacking has written
      permission.
    \item Social engineering tricks the person; pretexting invents a
      story and baiting offers a lure.
    \item Phishing is fake-message fraud for credentials or money;
      spoofing is the forged identity it may wear.
    \item Password attacks include guessing, dictionary, brute force,
      and shoulder surfing; long unique passwords plus two-step
      verification resist them.
    \item Malicious links and attachments carry the payload; verify the
      sender, type the address, and never share an OTP.
    \item Vulnerability $\rightarrow$ exploit $\rightarrow$ attack
      $\rightarrow$ payload is the chain to remember.
  \end{itemize}
\end{frame}
\end{document}
